\documentclass[conference]{IEEEtran}
\usepackage[T1]{fontenc}
\usepackage{cite}
\usepackage{amsmath,amssymb}
\usepackage{graphicx}
\usepackage{booktabs,array}
\usepackage{xcolor}
\usepackage{url}
\usepackage[hidelinks]{hyperref}
\hypersetup{pdftitle={When Operator Accuracy Does Not Transfer to the Image},
pdfauthor={Sally Mohlakola},pdfsubject={Learned interior transport for faceted gemstones}}
\graphicspath{{figures/}}
\newcommand{\dg}{^{\circ}}
\newcommand{\R}{\mathbb{R}}
\newcommand{\spp}{\,\mathrm{spp}}
\pagestyle{plain}
\raggedbottom

\begin{document}
\title{When Operator Accuracy Does Not Transfer to the Image: A Case Study of
a Learned Interior Transport Operator for Faceted Gemstones}
\author{\IEEEauthorblockN{Sally Mohlakola, [Supervisor~1] and [Supervisor~2]}
\IEEEauthorblockA{{[School], [University], [City], [Country]}\\
{[student email], [supervisor~1 email], [supervisor~2 email]}}}
\maketitle

\begin{abstract}
A diamond's appearance is made inside it: light refracts in, reflects between
facets and leaves in sharply defined directions. Neural appearance models
compress such internal light transport into a small network and work well for
cloth, where the structure is far smaller than a pixel. This study asks whether
the idea transfers to faceted gemstones, where the structure \emph{is} the
appearance. A radiance distribution map (RDM) adapted from the cloth method is
first shown to condition on the wrong variable, and is generalised to a
\emph{learned interior operator}: a network that takes where, in which direction
and at what wavelength light enters the stone, and predicts whether, through
which facet and in which direction it leaves. Three decoders were trained to
convergence on three diamond cuts with five seeds each (45 runs) and compared
with an exact spectral path tracer. The deterministic decoders predict exit
facets more accurately than the stochastic one on every cut, yet their images
agree with the reference \emph{less} (pixel correlation 0.050--0.184 against
0.122--0.272), in 29 of 30 paired seeds. Held-out paths locate the failure:
exit facets are predicted correctly for 45--59\% of paths meeting two or three
facets but 2--3\% of those meeting sixteen or more. Each learned decision is
three times faster than tracing, yet renders are 31\% slower and the neural
image settles at an error a single analytic render already beats.
\end{abstract}

\begin{IEEEkeywords}
neural rendering, light transport, gemstones, mixture density networks,
image evaluation, Monte Carlo rendering
\end{IEEEkeywords}

\begin{figure*}[t]
\centering
\includegraphics[width=\textwidth]{fig1_overview.png}
\caption{The learned interior operator. Camera rays meet the stone's first
surface, where Fresnel reflection or refraction is traced exactly. A refracted
ray's entry state (position, direction, normal, wavelength; 10 numbers) passes
through a shared two-layer encoder to an escape head, a facet head over the
mesh's triangles, and a 16-number embedding of the chosen facet. One of three
decoders then places the exit on that facet, and the exit ray returns to the
path tracer. The deep decoder alone also reads the entry state through its own
trunk.}
\label{fig:overview}
\end{figure*}

\section{Introduction}
\footnotetext{Source code available at: [GitHub repository].}
The appearance of a faceted diamond depends almost entirely on where light
leaves it. Light that enters the crown is refracted, reflected between facets
several times and refracted out again, and each facet acts as a mirror onto a
different part of the surroundings. Small changes in where light enters move
these reflections across facet edges, which is what produces brilliance, fire
and sparkle. Path tracing reproduces this exactly but repeats the whole
interior journey for every sample.

Neural appearance models offer a way to learn such journeys once and reuse
them. Soh and Montazeri~\cite{soh2024} compress the light scattered inside
cloth fibres into small networks trained from a radiance distribution map
(RDM): a table, gathered by tracing the true transport, of how light leaves for
each incoming direction. For cloth this is well founded: hundreds of fibres sit
inside one pixel, so light enters and leaves within a sub-pixel cross-section
and direction alone describes it. A diamond inverts this. It has a few dozen
facets, each covering thousands of pixels, and \emph{where} light enters
decides which facet it leaves through.

This study asks: \emph{can a compact learned operator reproduce the interior
light transport of a faceted gemstone, and if not, which of its errors explain
the difference in the image?} Three hypotheses make the question testable:
\begin{itemize}
\item \textbf{H1 (dissociation).} An operator can improve its exit-facet
accuracy without improving the image's agreement with the exact reference.
\item \textbf{H2 (localisation).} Its errors concentrate in long internal
routes and pavilion exits, not in facets that lack training data.
\item \textbf{H3 (cost).} A faster interior decision does not make the image
cheaper at equal quality.
\end{itemize}
All three are tested under one studio lighting rig on three procedural cuts;
the conclusions concern this setting rather than all learned representations of
specular media.

\section{Background}
\subsection{Light transport in a diamond}
When light meets a diamond's surface, part of it reflects and the rest
refracts into the stone. The split is given by the Fresnel equations and the
bend by Snell's law, $n_1\sin\theta_1=n_2\sin\theta_2$, with refractive index
$n\approx2.42$ for diamond. Inside, light meeting a facet beyond the critical
angle of about $24.4\dg$ is totally internally reflected, so it may bounce many
times before escaping. The index also varies with wavelength (dispersion),
which splits white light into the colours known as fire. A cut is designed so
that light entering the crown is reflected by the pavilion back out through the
crown; the result is a pattern of sharp reflections whose positions depend
sensitively on the viewpoint and on where light enters.

\subsection{Path tracing and Monte Carlo noise}
Physically based renderers estimate the light reaching each pixel by tracing
random light paths~\cite{kajiya1986,pharr2023}. Each pixel averages many
samples per pixel (spp); the estimate is unbiased but noisy, with variance that
falls in proportion to $1/N$ for $N$ samples. Two renders with different random
seeds therefore differ even when the method is exact, and any comparison
between images must account for this noise. For a gemstone, every camera path
that enters the stone must be traced through all of its internal bounces,
which is the work a learned operator aims to replace.

\subsection{Neural appearance models and the RDM}
A neural appearance model replaces part of the light transport with a small
network evaluated at render time~\cite{zeltner2024,kuznetsov2021}. In the cloth
model~\cite{soh2024}, the multiple scattering inside a yarn is gathered into an
RDM by firing rays at the fibre geometry and binning the light that leaves by
outgoing direction for each incoming direction; small networks are fitted to the
binned table and applied as a material. The present study starts from this
design and adapts it to a diamond.

\section{Related Work}
\subsection{Neural appearance and learned transport}
Neural materials compress measured or simulated appearance into compact
networks, from neural BRDFs with importance sampling~\cite{sztrajman2021} to
multi-resolution neural materials~\cite{kuznetsov2021} and real-time layered
models~\cite{zeltner2024}. These represent light that is reflected or scattered
near a surface. Closer to this study, Kallweit et al.~\cite{kallweit2017}
replace multiple scattering in clouds with a radiance-predicting network, and
Vicini et al.~\cite{vicini2019} learn where light leaves a translucent object
given where it enters, using a conditional variational autoencoder and local
shape descriptors. Both address diffusive media, whose exit distributions are
smooth. Learned densities also serve as importance-sampling
proposals~\cite{muller2019nis}. This study targets the specular case these
works exclude: a smooth dielectric whose exit is a discontinuous function of the
entry state.

\subsection{Rendering gemstones}
Gemstone rendering has been approached with explicit optics: Sun et
al.~\cite{sun2000} model the refraction, reflection and dispersion of diamonds,
and Guy and Soler~\cite{guy2004} render gemstones in real time by tracing a
bounded number of internal reflections. Spectral rendering with hero
wavelengths~\cite{wilkie2014} keeps dispersion affordable. These methods are
exact or deliberately truncated; none learns the interior transport.

\subsection{Representation limits of small networks}
Multilayer perceptrons are biased towards smooth functions~\cite{rahaman2019}.
Fourier features~\cite{tancik2020} and multiresolution hash
encodings~\cite{muller2022ngp} are standard remedies, and mixture density
networks~\cite{bishop1994} represent several possible outcomes for one input,
suiting an entry state that can escape by several internal routes. Moment-based
filtering such as LEAN mapping~\cite{olano2010} preserves variation that
averaging would otherwise erase. Each of these is tested or adapted below.

\subsection{Evaluating renders}
Images are commonly compared with SSIM~\cite{wang2004} or
FLIP~\cite{andersson2020}. Both are sensitive to Monte Carlo noise, which
matters when a learned method is compared with a noisy reference. This study
instead measures agreement against the reference's own noise level, corrects
correlations for noise in the manner of Spearman's
disattenuation~\cite{spearman1904}, and attaches block-bootstrap
intervals~\cite{efron1993,kunsch1989}.

\section{Preliminary Study: The Radiance Map}
\label{sec:rdm}
The project began as a direct adaptation of the cloth model: an RDM gathered on
a diamond mesh, compressed into small networks and applied as a material on the
stone's surface. Each step below removed one explanation for why this did not
produce a diamond.

\begin{figure*}[t]
\centering
\includegraphics[width=0.78\textwidth]{rdm_ladder.pdf}
\caption{RDM renders against gathering budget, after the 50--100 million rays
per material of the cloth model~\cite{soh2024}. For each cut the large panel
uses $2^{30}$ rays and the stacked panels $2^{26}$ to $2^{29}$; all use the same
studio rig and render settings. \emph{Placeholder layout: renders pending.}}
\label{fig:ladder}
\end{figure*}

\paragraph{Is it a sampling problem?} Soh and Montazeri gather 50--100 million
fully traced rays per yarn material until the map converges~\cite{soh2024}. To
give the RDM the same chance, it is gathered for each cut at five budgets from
$2^{26}$ rays (67 million) to $2^{30}$ (1.07 billion) and rendered under the
same rig (Fig.~\ref{fig:ladder}). If flat shading were sampling noise it would
sharpen along this ladder; if it persists while the bins converge, the limit is
what the map conditions on. \emph{[Ladder renders pending.]}

\paragraph{Averaging erases the detail} Each bin stores the mean of the paths
landing in it. In 99.8\% of well-populated bins the standard deviation exceeds
that mean (median ratio 2.53): more rays make the mean more precise but do not
make the distribution narrower. A moment-based modification after LEAN
mapping~\cite{olano2010}, which switches shading points on with probability
$p=1/(1+r^2)$ at gain $1/p$ for relative spread $r$, restores the \emph{amount}
of sparkle but places it arbitrarily: the result reads as speckle.

\paragraph{The aim offset fixes coverage} The first gathers aimed every ray at
the stone's centre, which caps the incidence angle each facet can receive: no
checkpoint before the fix filled more than 41 of its 64 incident-direction bins.
Firing each ray from a random point on a disc the size of the stone's
silhouette (the \emph{aim offset}), with directions stratified uniformly in
$\cos\theta$ within each bin, fills all 64. Where a ray is aimed matters as much
as its direction, a lesson carried into the operator's data (Sec.~\ref{sec:sampling}).

\paragraph{The missing variable is position} Asked to predict the exit
triangle, a table conditioned on entry direction, as the RDM is, scores 0.068
with 32 cells and only 0.098 with 2,048, against 0.047 for the most common
triangle; conditioning on the entry triangle instead reaches 0.192. Resolution
is not the limit; the conditioning variable is. The learned interior operator
follows: it keeps the RDM's idea of a learned interior but conditions on entry
position, direction, normal and wavelength and returns a single exit rather
than a binned average.

\section{Methodology}
\label{sec:method}
\subsection{Scene and reference}
All renders use one studio scene: a perspective camera with a $30\dg$ field of
view, a grey diffuse ground plane, textured area lights and two small intense
sources. Structured lighting is essential, because under uniform lighting every
exit direction returns the same radiance and errors become invisible. The
\emph{analytic reference} is a spectral path tracer in Mitsuba~3~\cite{mitsuba3}
that treats the stone as a smooth dispersive dielectric (Sellmeier index),
samples Fresnel reflection or refraction at every interaction and traces each
interior path until escape or a depth of 128. Absorption is not modelled.

\subsection{The learned interior operator}
The entry state $\mathbf{x}\in\R^{10}$ holds position $\mathbf{p}$ (scaled by
the mesh radius), interior direction $\boldsymbol\omega$, entry normal
$\mathbf{n}$ and wavelength $\lambda$. With $e$ the event that light escapes
before the depth cap, $f$ the exit triangle and $\mathbf{y}\in\R^4$ the exit
coordinates, the operator models
\begin{equation}
 p(e=1,f,\mathbf{y}\mid\mathbf{x})=
 s_\theta(\mathbf{x})\,\pi_\theta(f\mid\mathbf{x})\,
 q_\theta(\mathbf{y}\mid\mathbf{x},f),
\end{equation}
so non-escape is explicit rather than an artificial exit. Two coordinates are
barycentric logits and two are tangent slopes on triangle $f$:
\begin{equation}
 \boldsymbol\beta=\operatorname{softmax}(y_1,y_2,0),\quad
 \boldsymbol\omega_{\rm out}=
 \frac{\mathbf{n}_f+y_3\mathbf{t}_f+y_4\mathbf{b}_f}
 {\|\mathbf{n}_f+y_3\mathbf{t}_f+y_4\mathbf{b}_f\|},
\end{equation}
with exit point $\sum_j\beta_j\mathbf{v}_{fj}$, so every prediction lies on the
chosen triangle and points outward.

Fig.~\ref{fig:overview} shows the architecture. A two-layer, width-64 SiLU
encoder~\cite{elfwing2018} feeds an escape head and a facet head; a learned
16-number embedding of the chosen facet conditions the exit decoder. Three
decoders share everything else: the \emph{stochastic} decoder is a mixture of
four Gaussians~\cite{bishop1994} sampled at render time (17,637 parameters on
the round cut); the \emph{clone} regresses a single exit with a smooth-L1 loss,
after behavioural cloning~\cite{pomerleau1989} (15,557); and the \emph{deep
clone} uses three hidden layers and its own encoder (28,741). All three sample
the exit facet; they differ only in how the exit is placed on it. Training
minimises
\begin{equation}
 \mathcal L_i=\operatorname{BCE}(s_\theta(\mathbf{x}_i),e_i)
 +e_i\left[-\log\pi_\theta(f_i\mid\mathbf{x}_i)
 +\ell_{\rm coord}(\mathbf{y}_i)\right]
\end{equation}
with Adam~\cite{kingma2015}.

\subsection{Rendering with the operator}
The first surface is traced exactly in both modes; only the journey inside the
stone differs. When a camera path reaches the stone, reflection or refraction is
sampled from the Fresnel coefficient. A reflected path continues outside as
usual. A refracted path forms its entry state and joins a batch of up to 512
waiting entries, so the network is queried efficiently. For each entry the
operator samples whether light escapes, the exit facet and the exit
coordinates; the path then continues from the decoded exit as an ordinary scene
ray. Non-escaping paths end, and exits that would immediately re-enter the stone
are discarded. In analytic mode the same entry is instead traced through every
internal bounce.

\section{Experiments}
\begin{table}[t]
\caption{Experimental parameters}
\label{tab:params}
\centering\small
\setlength{\tabcolsep}{3pt}
\begin{tabular}{@{}lp{46mm}@{}}
\toprule
Parameter & Value \\
\midrule
Cuts (triangles / facets) & round 64/57, pear 62/33, step 112/50 \\
Training pool per cut & 1,048,576 paths from 65,536 entry states: half camera-matched, half uniform with the aim offset; split 80/20 by entry state \\
Held-out test pool per cut & 262,144 paths from 16,384 entry states, same mix, separate seeds \\
Interior depth cap & 128 \\
Training & Adam, learning rate $10^{-3}$, batch 256, clip 10 \\
Stopping & best validation loss improving by $\le10^{-3}$ for 5 epochs (stopped after 31--93) \\
Training seeds & 23, 24, 25, 26, 27 \\
Renders & $512\times512$, box filter, $8\times32\spp$ per mode (halves of $128\spp$) \\
Operator batch & 512 entry states \\
Hardware & 14-core Linux nodes of a Slurm cluster \\
Software & Mitsuba 3.9.1, Dr.Jit 1.5.0, PyTorch 2.13, Python 3.10 \\
\bottomrule
\end{tabular}
\end{table}

\begin{figure*}[t]
\centering
\includegraphics[width=\textwidth]{gia_three_heads.png}
\caption{Round-cut renders from E1 (seed 23): stochastic, clone and deep decoder
beside the analytic reference ($256\spp$). Deterministic placement restores sharp
features but not their locations. Panel scores are that run's pixel
correlation; Table~\ref{tab:images} gives five-seed means.}
\label{fig:round}
\end{figure*}

\begin{table*}[t]
\caption{Image agreement in E1: mean $\pm$ standard deviation over five training
seeds. Each render is $8\times32\spp$; pixel, disattenuated and interior-only
correlations compare $128\spp$ halves, facet correlation and noise-free contrast
use the full average. Floors: the analytic reference against its own other half.}
\label{tab:images}
\centering\small
\begin{tabular}{@{}llrrrrrr@{}}
\toprule
Cut & Decoder & Pixel corr. $\uparrow$ & Disattenuated $\uparrow$ & Interior only $\uparrow$ & Facet corr. $\uparrow$ & Energy error & Contrast \\
\midrule
Round & Stochastic & \textbf{0.122} $\pm$ 0.015 & \textbf{0.154} $\pm$ 0.024 & \textbf{0.057} $\pm$ 0.014 & \textbf{0.260} $\pm$ 0.064 & $+8\%$ $\pm$ 2 & 0.85 $\pm$ 0.09 \\
 & Clone & 0.050 $\pm$ 0.025 & 0.053 $\pm$ 0.027 & 0.021 $\pm$ 0.017 & 0.250 $\pm$ 0.122 & $+21\%$ $\pm$ 18 & 2.22 $\pm$ 0.54 \\
 & Deep decoder & 0.060 $\pm$ 0.037 & 0.063 $\pm$ 0.038 & 0.032 $\pm$ 0.037 & 0.120 $\pm$ 0.076 & $+13\%$ $\pm$ 12 & 2.69 $\pm$ 0.71 \\
 & \emph{floor / reference} & \emph{0.974} & & & \emph{0.9999} & & \emph{3.72} \\
\midrule
Pear & Stochastic & \textbf{0.172} $\pm$ 0.014 & \textbf{0.281} $\pm$ 0.020 & \textbf{0.103} $\pm$ 0.016 & \textbf{0.690} $\pm$ 0.072 & $+19\%$ $\pm$ 2 & 0.48 $\pm$ 0.02 \\
 & Clone & 0.074 $\pm$ 0.038 & 0.081 $\pm$ 0.042 & 0.034 $\pm$ 0.042 & 0.476 $\pm$ 0.214 & $+46\%$ $\pm$ 17 & 1.74 $\pm$ 0.26 \\
 & Deep decoder & 0.092 $\pm$ 0.019 & 0.103 $\pm$ 0.022 & 0.050 $\pm$ 0.018 & 0.451 $\pm$ 0.138 & $+40\%$ $\pm$ 5 & 1.46 $\pm$ 0.22 \\
 & \emph{floor / reference} & \emph{0.941} & & & \emph{0.9999} & & \emph{1.44} \\
\midrule
Step & Stochastic & \textbf{0.272} $\pm$ 0.018 & \textbf{0.356} $\pm$ 0.021 & \textbf{0.197} $\pm$ 0.020 & \textbf{0.857} $\pm$ 0.023 & $+16\%$ $\pm$ 3 & 0.59 $\pm$ 0.02 \\
 & Clone & 0.112 $\pm$ 0.036 & 0.121 $\pm$ 0.040 & 0.058 $\pm$ 0.027 & 0.713 $\pm$ 0.087 & $+29\%$ $\pm$ 9 & 1.53 $\pm$ 0.36 \\
 & Deep decoder & 0.184 $\pm$ 0.026 & 0.203 $\pm$ 0.031 & 0.124 $\pm$ 0.026 & 0.852 $\pm$ 0.031 & $+15\%$ $\pm$ 4 & 1.09 $\pm$ 0.14 \\
 & \emph{floor / reference} & \emph{0.975} & & & \emph{0.9999} & & \emph{2.35} \\
\bottomrule
\end{tabular}
\end{table*}

\begin{figure*}[t]
\centering
\includegraphics[width=\textwidth]{operator_breakdown.pdf}
\caption{E3, round cut, stochastic decoder, seed 23: 261,335 escaped held-out
paths. (a) Exit-facet top-1 accuracy by the number of facets a path meets inside
the stone; the second tick line is each group's share of paths. (b) Exit-direction
error when the true facet is supplied: median and 90th percentile. (c) Share of
exits through each region that the operator also sends there; bar: whole region,
dot: one triangle with at least 20 test exits.}
\label{fig:breakdown}
\end{figure*}

\subsection{Experimental design}
Six experiments test the hypotheses.
\begin{itemize}
\item \textbf{E1 Decoder comparison} (H1): the stochastic, clone and deep
decoders on each of the three cuts, each trained with five seeds (45 runs), all
trained to convergence on the same pool per cut, tested on an independent pool
and rendered against one analytic reference per cut.
\item \textbf{E2 Diagnostics} (H1): \emph{true-facet} renders, in which the
exact tracer supplies each path's exit facet to the seed-23 models; and,
from earlier single-seed experiments on a camera-matched pool, lookup baselines,
a nearest-neighbour oracle and an input ablation.
\item \textbf{E3 Breakdown} (H2): E1's held-out paths grouped by the number of
facets met inside the stone and by exit region.
\item \textbf{E4 Cost} (H3): render and per-decision timing at equal time and
equal quality, and the network's cost against batch size.
\item \textbf{E5 Entry sampling}: camera-matched against uniform aim-offset
training entries, on a 60-frame pear tumble.
\item \textbf{E6 Transfer}: an operator trained on one cut, re-pointed at
another without retraining: round (GIA) to pear, and pear to step cut.
\end{itemize}

\subsection{Baselines}
\begin{itemize}
\item \textbf{Analytic reference}: the exact path tracer at the same settings.
\item \textbf{Lookup tables}: the most common exit triangle, and tables keyed on
entry direction (the RDM's conditioning) or entry triangle.
\item \textbf{Nearest-neighbour oracle}: the exits of the closest training
entry states, a non-parametric test of whether the exit is predictable at all.
\end{itemize}

\begin{table}[t]
\caption{Round-cut exit-triangle prediction on held-out paths. Upper rows:
earlier single-seed diagnostics on a camera-matched pool; lower rows: E1, mean
$\pm$ standard deviation over five seeds on the mixed test pool.}
\label{tab:baselines}
\centering\small
\begin{tabular}{@{}lr@{}}
\toprule
Predictor & Top-1 accuracy \\
\midrule
Most common triangle & 0.047 \\
Entry-direction table (32 / 2,048 cells) & 0.068 / 0.098 \\
Entry-triangle table & 0.192 \\
Nearest neighbour / majority of eight & 0.287 / 0.343 \\
\midrule
Stochastic decoder & 0.358 $\pm$ 0.005 \\
Clone & 0.386 $\pm$ 0.005 \\
Deep decoder & 0.386 $\pm$ 0.003 \\
\bottomrule
\end{tabular}
\end{table}

\begin{figure}[t]
\centering
\includegraphics[width=\columnwidth]{pear_heads.png}\\[2pt]
\includegraphics[width=\columnwidth]{step_heads.png}
\caption{Pear (top) and step cut (bottom) from E1 (seed 23): stochastic and
deep-decoder operators beside their references ($256\spp$). Broad shape and
brightness survive on both cuts; the sharp internal reflections do not.}
\label{fig:cuts}
\end{figure}
\begin{figure}[t]
\centering
\includegraphics[width=0.92\columnwidth]{oracle_facet.png}
\caption{E2, round cut, seed 23: the clone, the clone supplied with the traced
exit facet, and the reference ($256\spp$). The correct facet restores the star
pattern and contrast but not placement.}
\label{fig:oracle}
\end{figure}

\begin{table}[t]
\caption{E2: supplying the true exit facet to the seed-23 models. Pixel
correlation and noise-free contrast, without and with the true facet. Step cut
pending.}
\label{tab:truefacet}
\centering\small
\setlength{\tabcolsep}{3.5pt}
\begin{tabular}{@{}llrrrr@{}}
\toprule
& & \multicolumn{2}{c}{Pixel corr.} & \multicolumn{2}{c}{Contrast (ref.)} \\
\cmidrule(lr){3-4}\cmidrule(l){5-6}
Cut & Decoder & learned & true facet & learned & true facet \\
\midrule
Round & Stochastic & 0.127 & 0.151 & 0.79 & 0.98 (3.72) \\
 & Clone & 0.077 & 0.048 & 1.73 & 2.73 (3.72) \\
 & Deep & 0.057 & 0.077 & 1.53 & 2.50 (3.72) \\
\midrule
Pear & Stochastic & 0.188 & 0.310 & 0.50 & 0.84 (1.44) \\
 & Clone & 0.121 & 0.127 & 1.94 & 4.90 (1.44) \\
 & Deep & 0.107 & 0.107 & 1.41 & 3.38 (1.44) \\
\bottomrule
\end{tabular}
\end{table}

\subsection{Evaluation metrics}
Metrics use the pixels whose camera ray hits the stone, on linear luminance.
Eight seeds per mode are split into two independent halves; the reference
compared with its own other half gives the \emph{noise floor}, the agreement
sampling noise alone allows.
\begin{itemize}
\item \textbf{Pixel correlation} $\rho(N,A)$ between neural and reference
halves: is light in the right place?
\item \textbf{Disattenuated correlation} $\rho/\sqrt{r_Nr_A}$, dividing by
each image's half-split reliability $r$: agreement with noise removed.
\item \textbf{Facet correlation} across visible facet means: is coarse
brightness right?
\item \textbf{Energy error} $\epsilon_E=\sum N_p/\sum A_p-1$: is the total right?
\item \textbf{Contrast}, standard deviation over mean (noise-free version from
the covariance of halves): is the image as sharp as the reference?
\item \textbf{Interior-only} versions of each, using only light that passed
through the operator (or, for the reference, the traced interior).
\item \textbf{Operator metrics} on held-out paths: exit-facet top-1 accuracy,
and the exit-direction angle when the true facet is supplied.
\end{itemize}
Decoders are compared seed by seed: each deterministic run is paired with the
stochastic run of the same seed, and a difference counts as clear when its
95\% block-bootstrap interval excludes zero.

\subsection{Data collection}
Each of the 45 runs writes, under one folder: the exact command and git commit,
the software and hardware, hashes of the mesh, training data and checkpoint,
per-epoch training and validation losses, per-seed render timing, every
evaluation result as JSON and CSV, and its figures. A final step aggregates the
five seeds of every cut and decoder and computes the paired differences. The
runs were executed as Slurm array jobs; every job skips finished work, so
interrupted jobs resume without repeating it.

\section{Results and Analysis}
\begin{table}[t]
\caption{E4 cost, round cut, stochastic decoder, seed 23: $512\times512$,
32 spp per render, mean of eight seeds. Error is relative RMSE against the
converged analytic image.}
\label{tab:cost}
\centering\small
\setlength{\tabcolsep}{4pt}
\begin{tabular}{@{}lrr@{}}
\toprule
& Analytic & Neural \\
\midrule
Render time (s) & \textbf{264} & 345 \\
Time per interior decision ($\mu$s) & 31 & \textbf{10.5} \\
Interior share of render time & 50\% & 13\% \\
Error at equal time (2,760 s) & \textbf{0.36} & 3.72 \\
Lowest reachable error & $\to 0$ & 3.69 \\
Time to reach the other's best & \textbf{26 s} & never \\
Training time & -- & 485 s \\
Model size & 1 KB mesh & 70 KB \\
\bottomrule
\end{tabular}
\end{table}

\begin{figure}[t]
\centering
\includegraphics[width=0.84\columnwidth]{timing.pdf}
\caption{E4, where the time goes. (a) Render time split into interior decisions
and everything else, with the seed range. (b) Cost of one interior decision
against batch size: the network alone (curve), the analytic trace in the
renderer (dashed) and the network in the renderer at its mean batch (diamond).}
\label{fig:timing}
\end{figure}

\begin{figure}[t]
\centering
\includegraphics[width=0.92\columnwidth]{sampling_contrast.pdf}
\caption{E5, entry sampling on the pear. Top: where training entries land under
(a) uniform directions with the aim offset and (b) camera rays from a full
orbit. Bottom: a pavilion-facing tumble pose (frame 53 of 60, $256\spp$)
rendered by operators trained on (b) alone and on (a) and (b) together.}
\label{fig:sampling}
\end{figure}

\subsection{The operator learns (E1, E2)}

\paragraph{Learning beats every lookup} All three decoders predict the round
cut's exit triangle better than every lookup baseline and the nearest-neighbour
oracle (Table~\ref{tab:baselines}), with a spread of at most 0.005 between seeds.
Training converged after 31--93 epochs, so the networks are not under-trained.

\paragraph{Position carries the information} In the earlier ablation, removing
position from the input dropped top-1 accuracy from 0.404 to 0.273, and an
operator given only entry direction, the RDM's conditioning, fell to 0.117.
Removing wavelength changed it by 0.0003.

\subsection{Sharper is not more faithful (E1, H1)}

\paragraph{A more accurate operator agrees less} On every cut the clone and deep
decoders predict exit facets more accurately than the stochastic decoder on the
held-out pool (round 0.386 against 0.358, pear 0.31 against 0.27, step 0.29
against 0.24), and their images look more like a diamond
(Figs.~\ref{fig:round} and~\ref{fig:cuts}). Yet their pixel correlation is lower on every cut
(Table~\ref{tab:images}). Paired by seed, the deterministic decoder is below
the stochastic one in 29 of 30 pairs, and clearly so in 26. This supports H1.

\paragraph{The gap is not noise} Removing sampling noise (disattenuated
correlation) leaves the ordering unchanged, and every score stays far below the
noise floors of 0.94--0.98. Light that actually passes through the network
agrees least: interior-only correlations are about half the whole-stone values.

\paragraph{The stochastic image is right on average, wrong in detail} The
stochastic decoder keeps the best agreement but a quarter to a third of the
reference's contrast (0.85 against 3.72 on the round cut): its exits spread so
widely that each pixel averages bright and dark parts of the lighting.

\paragraph{Coarse structure can survive} Averaged over facets, the step cut's
deep decoder matches the stochastic one (0.852 against 0.857) while its pixel
correlation is lower: which facets are bright can be right while the sparkle
within them is wrong.

\subsection{Choosing the facet is not enough (E2)}

\paragraph{The right facet restores the look} Supplying the true exit facet
raises the deterministic decoders' contrast towards or beyond the reference's
(Table~\ref{tab:truefacet}, Fig.~\ref{fig:oracle}): on the pear the clone
reaches 4.90 against 1.44.

\paragraph{The right facet does not fix agreement} Their pixel correlation
stays between 0.05 and 0.13, and the round clone's even falls (0.077 to 0.048),
because the exit is still misplaced within the facet. Only the stochastic decoder gains clearly (pear 0.188 to 0.310), as its
broad placement keeps overlap once the facet is right. Facet choice and
within-facet placement are therefore separate failures.

\subsection{Where the operator fails (E3, H2)}

\paragraph{Accuracy collapses on long paths} On all three cuts, paths meeting
two or three facets are predicted correctly 45--59\% of the time and leave
within a median $9.5\dg$--$11\dg$ of the traced direction, but paths meeting
sixteen or more are predicted correctly 2--3\% of the time and miss by
$44\dg$--$48\dg$ (Fig.~\ref{fig:breakdown}a,b). This part of H2 holds on every
cut and every seed.

\paragraph{The pavilion is not the problem once it is covered} H2 also
predicted failure at pavilion exits. Trained on the mixed pool, pavilion exits
are predicted as well as crown exits (seed 23: 29\%, 27\% and 28\% against 32\%,
32\% and 16\% on the round, pear and step cuts; Fig.~\ref{fig:breakdown}c). An
earlier camera-only model predicted only 6\% of pavilion exits correctly, so that
failure was a coverage effect (E5), not a property of the pavilion. H2 is
supported for route length but not for exit region.

\subsection{Faster per decision, slower per image (E4, H3)}
\label{sec:cost}

\paragraph{Each decision is faster} One learned interior decision costs
10.5\,$\mu$s against 31\,$\mu$s to trace it, three times faster, and the
interior's share of the render drops from 50\% to 13\% (Table~\ref{tab:cost},
Fig.~\ref{fig:timing}a).

\paragraph{The whole render is slower} A neural render takes 31\% longer: the
time outside the interior grows from 132\,s to 301\,s, spent gathering entry
states, preparing inputs and decoding exits. Batch size limits the saving
(Fig.~\ref{fig:timing}b): one query alone costs about 0.4\,ms, the network only
overtakes tracing beyond about 17 queries per batch, and although the renderer
asks for 512, its batches average 83 as paths finish.

\paragraph{No budget reaches analytic quality} The neural error is a fixed bias
plus noise and levels off at 3.69, whereas the analytic image reaches that error
after about 26\,s. At equal time the analytic error is 10 times lower on the
round cut and 7--10 times lower across the cuts. This supports H3.

\subsection{Entry sampling decides coverage (E5)}
\label{sec:sampling}

\paragraph{Camera sampling misses the pavilion} Camera-matched entries match
what the renderer asks, but a camera above the stone sees only the crown, so no
entries land on the pavilion; uniform sampling with the RDM's aim offset puts
54\% there (Fig.~\ref{fig:sampling}). Trained on camera entries alone, the
operator lets only 46\% of queries escape when the pavilion faces the camera,
and the pavilion renders nearly black.

\paragraph{Coverage repairs escape, not placement} Adding uniform aim-offset
entries restores escape to 100\% and lights the pavilion. In earlier tests the
correlation at that pose stayed at 0.018 against a floor of 0.951. This is why
every E1 pool mixes the two strategies.

\subsection{Transfer across cuts (E6)}
\label{sec:transfer}
\begin{figure}[t]
\centering
\includegraphics[width=\columnwidth]{transfer.pdf}
\caption{E6: an operator trained on one cut and re-pointed at another without
retraining, beside the stochastic operator trained on the target cut and the
analytic reference ($256\spp$).}
\label{fig:transfer}
\end{figure}

\paragraph{What transfer requires} The operator's facet head is a categorical
over the triangles of its training mesh, so each target triangle is given the
identity of the source triangle whose outward normal points most nearly the same
way, and the head's weights and facet embeddings are permuted to match; nothing
is retrained. Round to pear (64 to 62 triangles) drops two source classes; pear
to step cut (62 to 112) lets 50 surplus triangles share the probability of their
best-aligned source triangle. Matching by normal ignores position, and the
network keeps the internal routes of the source stone.

\paragraph{Round to pear partly transfers} The seed-23 round operator on the pear
(normals matched within a median $7.8\dg$, mean $30.9\dg$) keeps much of the
coarse facet brightness (facet correlation 0.553 against 0.690 for the pear's own
operator), but its pixel correlation falls from 0.172 to 0.070, its
interior-only correlation to $-0.005$, and it adds 44\% too much energy.

\paragraph{Pear to step cut does not} Across families the transfer fails
outright: pixel correlation $-0.018$, interior-only $-0.040$ and facet
correlation 0.001, against 0.272, 0.197 and 0.857 for the step cut's own
operator. Learned routes do not carry over to a different cut family.
\emph{[Deep-decoder transfers pending.]}

\section{Discussion}
\subsection{Why accuracy does not reach the image}
For entry state $\mathbf{x}$ and exit state $z$, the light carried back is
\begin{equation}
 L(\mathbf{x})=\int L_{\rm ext}(z)\,p(z\mid\mathbf{x})\,\mathrm{d}z.
\end{equation}
The training loss weights every exit equally, but the image weights each exit
by the radiance $L_{\rm ext}$ it reaches: a small angular error at the edge of a
bright lamp matters greatly, a large one between dark directions hardly at all.
Facet accuracy and mean angular error are therefore incomplete proxies for the
image. A broad exit distribution spreads energy over correct and incorrect
regions and keeps some overlap; a sharp but misplaced one loses it, which is why
the sharper decoders score lower in E1 and gain least from the true facet in E2.
This explanation is consistent with the results but untested; a
radiance-weighted objective would test it.

E3 adds that accuracy follows route length on every cut once coverage is fixed.
A long route changes facet whenever the path crosses any internal facet edge, so
its exit varies rapidly with the entry state, and a small smooth network blends
across such changes~\cite{rahaman2019}. Local encodings~\cite{muller2022ngp} are
a natural next test.

\subsection{Limitations}
The study uses one lighting rig, one view for the main comparison, simplified
meshes and a facet head tied to its mesh. Five training seeds bound the
seed-to-seed spread but leave the intervals wide. The lookup baselines,
nearest-neighbour oracle and ablation come from earlier single-seed runs on a
camera-matched pool. The step-cut true-facet renders, the deep-decoder transfers
and the RDM ladder are pending, and ground caustics are not evaluated.

\section{Conclusion}
This study adapted a neural appearance model from cloth to faceted gemstones
and generalised its direction-conditioned radiance map into a position-aware
learned interior operator, evaluated in 45 converged runs.

The operator learns: all three decoders predict exit facets better than every
lookup baseline, with a seed spread of at most 0.007.

Accuracy does not transfer to the image (H1). The more accurate deterministic
decoders render sharper diamonds that agree with the reference less on every cut
(0.050--0.184 against 0.122--0.272), in 29 of 30 paired seeds, and the true exit
facet restores their contrast but not their agreement.

The failure follows route length (H2, in part). Paths meeting two or three
facets are predicted correctly about half the time, those meeting sixteen or
more almost never; pavilion exits fail only when the pavilion is missing from
training.

The operator does not save time (H3). Each decision is three times faster, yet
renders are 31\% slower and the image settles at an error the analytic renderer
reaches in under half a minute.

Future work should test radiance-weighted training and local encodings on the
long-route cases, and repeat the study under other lighting.

\begin{thebibliography}{99}
\bibitem{soh2024} G. Y. Soh and Z. Montazeri, ``Neural appearance model for cloth
rendering,'' \emph{Computer Graphics Forum}, vol. 43, no. 4, 2024,
doi: 10.1111/cgf.15156.
\bibitem{kajiya1986} J. T. Kajiya, ``The rendering equation,'' \emph{Computer
Graphics (Proc. SIGGRAPH)}, vol. 20, no. 4, pp. 143--150, 1986.
\bibitem{pharr2023} M. Pharr, W. Jakob, and G. Humphreys,
\emph{Physically Based Rendering: From Theory to Implementation}, 4th ed.
MIT Press, 2023.
\bibitem{zeltner2024} T. Zeltner \emph{et al.}, ``Real-time neural appearance
models,'' \emph{ACM Trans. Graph.}, vol. 43, no. 3, 2024.
\bibitem{kuznetsov2021} A. Kuznetsov \emph{et al.}, ``NeuMIP: Multi-resolution
neural materials,'' \emph{ACM Trans. Graph.}, vol. 40, no. 4, 2021.
\bibitem{sztrajman2021} A. Sztrajman, G. Rainer, T. Ritschel, and T. Weyrich,
``Neural BRDF representation and importance sampling,'' \emph{Computer Graphics
Forum}, vol. 40, no. 6, 2021.
\bibitem{kallweit2017} S. Kallweit, T. M\"uller, B. McWilliams, M. Gross, and
J. Nov\'ak, ``Deep scattering: Rendering atmospheric clouds with
radiance-predicting neural networks,'' \emph{ACM Trans. Graph.}, vol. 36, no. 6,
2017.
\bibitem{vicini2019} D. Vicini, V. Koltun, and W. Jakob, ``A learned shape-adaptive
subsurface scattering model,'' \emph{ACM Trans. Graph.}, vol. 38, no. 4,
Art. 126, 2019.
\bibitem{muller2019nis} T. M\"uller, B. McWilliams, F. Rousselle, M. Gross, and
J. Nov\'ak, ``Neural importance sampling,'' \emph{ACM Trans. Graph.}, vol. 38,
no. 5, 2019.
\bibitem{sun2000} Y. Sun, F. D. Fracchia, and M. S. Drew, ``Rendering diamonds,''
in \emph{Proc. 11th Western Computer Graphics Symposium}, 2000.
\bibitem{guy2004} S. Guy and C. Soler, ``Graphics gems revisited: Fast and
physically-based rendering of gemstones,'' \emph{ACM Trans. Graph.}, vol. 23,
no. 3, pp. 231--238, 2004.
\bibitem{wilkie2014} A. Wilkie, S. Nawaz, M. Droske, A. Weidlich, and J. Hanika,
``Hero wavelength spectral sampling,'' \emph{Computer Graphics Forum}, vol. 33,
no. 4, pp. 123--131, 2014.
\bibitem{rahaman2019} N. Rahaman \emph{et al.}, ``On the spectral bias of neural
networks,'' in \emph{Proc. Int. Conf. Machine Learning}, 2019, pp. 5301--5310.
\bibitem{tancik2020} M. Tancik \emph{et al.}, ``Fourier features let networks learn
high frequency functions in low dimensional domains,'' in \emph{Advances in
Neural Information Processing Systems}, vol. 33, 2020, pp. 7537--7547.
\bibitem{muller2022ngp} T. M\"uller, A. Evans, C. Schied, and A. Keller,
``Instant neural graphics primitives with a multiresolution hash encoding,''
\emph{ACM Trans. Graph.}, vol. 41, no. 4, Art. 102, 2022.
\bibitem{bishop1994} C. M. Bishop, ``Mixture density networks,'' Aston University,
Tech. Rep. NCRG/94/004, 1994.
\bibitem{olano2010} M. Olano and D. Baker, ``LEAN mapping,'' in \emph{Proc. ACM
SIGGRAPH Symp. Interactive 3D Graphics and Games}, 2010, pp. 181--188.
\bibitem{wang2004} Z. Wang, A. C. Bovik, H. R. Sheikh, and E. P. Simoncelli,
``Image quality assessment: From error visibility to structural similarity,''
\emph{IEEE Trans. Image Process.}, vol. 13, no. 4, pp. 600--612, 2004.
\bibitem{andersson2020} P. Andersson \emph{et al.}, ``FLIP: A difference evaluator
for alternating images,'' \emph{Proc. ACM Comput. Graph. Interact. Tech.},
vol. 3, no. 2, Art. 15, 2020.
\bibitem{spearman1904} C. Spearman, ``The proof and measurement of association
between two things,'' \emph{American Journal of Psychology}, vol. 15, no. 1,
pp. 72--101, 1904.
\bibitem{efron1993} B. Efron and R. J. Tibshirani, \emph{An Introduction to the
Bootstrap}. Chapman \& Hall, 1993.
\bibitem{kunsch1989} H. R. K\"unsch, ``The jackknife and the bootstrap for general
stationary observations,'' \emph{Annals of Statistics}, vol. 17, no. 3,
pp. 1217--1241, 1989.
\bibitem{mitsuba3} W. Jakob \emph{et al.}, ``Mitsuba 3 renderer,'' 2022.
[Online]. Available: \url{https://mitsuba-renderer.org}
\bibitem{elfwing2018} S. Elfwing, E. Uchibe, and K. Doya, ``Sigmoid-weighted
linear units for neural network function approximation in reinforcement
learning,'' \emph{Neural Networks}, vol. 107, pp. 3--11, 2018.
\bibitem{pomerleau1989} D. A. Pomerleau, ``ALVINN: An autonomous land vehicle in
a neural network,'' in \emph{Advances in Neural Information Processing
Systems}, 1989.
\bibitem{kingma2015} D. P. Kingma and J. Ba, ``Adam: A method for stochastic
optimization,'' in \emph{Proc. Int. Conf. Learning Representations}, 2015.
\end{thebibliography}
\end{document}
